\documentclass[11pt]{article}

\usepackage[margin=1in]{geometry}
\usepackage{amsmath}
\usepackage{amssymb}
\usepackage{xcolor}
\usepackage[most]{tcolorbox}
\usepackage{enumitem}
\usepackage{titlesec}
\usepackage{hyperref}

% Links map to the structural header blue.
% bookmarks=false avoids the spaced-jobname .out failure on Windows/tectonic.
\hypersetup{colorlinks=true, linkcolor=headerblue, urlcolor=headerblue, bookmarks=false}

\tcbuselibrary{skins}

% 1.1 Palette
\definecolor{headerblue}{HTML}{1C569E}
\definecolor{accentgreen}{HTML}{2E7D32}
\definecolor{accentorange}{HTML}{E27A1E}
\definecolor{workpurple}{HTML}{A020F0}
\definecolor{warnred}{HTML}{C0392B}

% 1.2 Subsection Typography
\titleformat{\subsection}
  {\normalfont\large\bfseries\color{headerblue}}{\thesubsection}{1em}{}

% 2. Table of Contents Depth
\setcounter{secnumdepth}{2}
\setcounter{tocdepth}{2}

% 3. Box System (Unbreakable)
\newtcolorbox{topicsbox}{
    enhanced,
    colback=headerblue!6!white, colframe=headerblue,
    coltitle=white, fonttitle=\bfseries,
    title={Non-Zero Conditions at a Glance: Core Sub-Topics},
    boxrule=0.6pt, arc=2mm,
    left=3mm, right=3mm, top=2mm, bottom=2mm
}

\newtcolorbox{formulabox}[1][Master Formula]{
    enhanced,
    colback=accentorange!8!white, colframe=accentorange,
    coltitle=white, fonttitle=\bfseries,
    title={#1},
    boxrule=0.6pt, arc=1mm,
    left=3mm, right=3mm, top=2mm, bottom=2mm
}

\newtcolorbox{methodbox}[1][Method]{
    enhanced,
    colback=accentgreen!6!white, colframe=accentgreen,
    coltitle=white, fonttitle=\bfseries,
    title={#1},
    boxrule=0.6pt, arc=1mm,
    left=3mm, right=3mm, top=2mm, bottom=2mm
}

\newtcolorbox{examplebox}[1][Worked Example]{
    enhanced,
    colback=workpurple!4!white, colframe=workpurple,
    coltitle=white, fonttitle=\bfseries,
    title={#1},
    boxrule=0.4pt, sharp corners,
    borderline west={3pt}{0pt}{workpurple},
    left=5mm, right=3mm, top=2mm, bottom=2mm
}

\newtcolorbox{conceptbox}[1][Key Takeaway]{
    enhanced,
    colback=accentorange!5!white, colframe=accentorange,
    coltitle=white, fonttitle=\bfseries,
    title={#1},
    boxrule=0.3pt, arc=1mm,
    left=3mm, right=3mm, top=2mm, bottom=2mm
}

\newtcolorbox{warningbox}[1][Common Pitfall]{
    enhanced,
    colback=red!3!white, colframe=warnred,
    coltitle=white, fonttitle=\bfseries,
    title={#1},
    boxrule=0.4pt, sharp corners,
    borderline west={3pt}{0pt}{warnred},
    left=5mm, right=3mm, top=2mm, bottom=2mm
}

\begin{document}

\begin{center}
    {\Large\bfseries\color{headerblue} Laplace Transforms with Non-Zero\\[2pt] Initial Conditions and Boundary Values}\\[6pt]
    \large A Self-Study Reference \& Master Guide
\end{center}
\vspace{0.4cm}

\tableofcontents
\vspace{0.8cm}

%======================================================================
\phantomsection
\addcontentsline{toc}{section}{Part I --- Theory \& Worked Examples}
\section*{Part I --- Theory \& Worked Examples}
\setcounter{section}{2}
\setcounter{subsection}{9}

\begin{topicsbox}
\begin{itemize}[leftmargin=*, itemsep=2pt]
  \item \textbf{The Hard-Wired Lower Bound (2.10):} the Laplace integral starts at $t = 0$, and nothing can move it. The derivative rule $\mathcal{L}\{y'\} = sY(s) - y(0)$ feeds on the value \emph{at zero} --- so a condition given anywhere else cannot enter the transform directly.
  \item \textbf{The ``Clock Reset'' Substitution (2.11):} for an IVP whose data live at $t = a \ne 0$, slide the time axis with $\tau = t - a$. The shifted problem has its conditions at $\tau = 0$, exactly where Laplace can use them. Solve in $\tau$, then slide back.
  \item \textbf{Parameterizing Boundary Value Problems (2.12):} when one condition sits at $t = 0$ and another at $t = b$, the missing initial slope is unknown. Carry it as a symbol $c$, solve the ODE in terms of $c$, then pin $c$ down with the far boundary condition.
\end{itemize}
\end{topicsbox}

\subsection{The $t=0$ Trap --- Why the Standard Rule Needs Data at Zero}

Every Laplace-based ODE solve rests on the derivative rule
\[ \mathcal{L}\{y'(t)\} = sY(s) - y(0), \qquad \mathcal{L}\{y''(t)\} = s^2 Y(s) - s\,y(0) - y'(0). \]
These constants $y(0)$ and $y'(0)$ are not a stylistic choice --- they fall directly out of integration by parts on the \emph{defining} integral, whose lower limit is permanently fixed at zero:
\[ \mathcal{L}\{f(t)\} = \int_{0}^{\infty} e^{-st}\, f(t)\, dt. \]
The boundary term from integration by parts is evaluated at that lower limit, which is why the rule asks for the value \emph{at $t = 0$} and nowhere else.

\begin{warningbox}[The $t=0$ Trap]
The derivative rule $\mathcal{L}\{y'\} = sY(s) - y(0)$ is \textbf{hard-wired to consume the value at zero}. The lower bound of the Laplace integral is fixed at $0$; it is part of the definition and cannot be slid to another point.

\smallskip
So if you are handed a condition such as $y(2) = 5$ --- a value at $t = 2$ --- there is \textbf{no slot} for it in the transform. Writing ``$sY(s) - y(2)$'' is meaningless: the rule does not have a $y(2)$ term, and inventing one corrupts the algebra. The transform only ever knows $y(0)$ and $y'(0)$.

\smallskip
\textbf{The fix depends on what you were given:}
\begin{itemize}[leftmargin=*, itemsep=1pt]
  \item \emph{All} conditions sit at one common point $t = a \ne 0$ (a shifted IVP) $\Rightarrow$ slide the clock so that point becomes $\tau = 0$ (Section 2.11).
  \item Conditions are \emph{split} between $t = 0$ and some $t = b$ (a boundary value problem) $\Rightarrow$ parameterize the unknown zero-data as a constant $c$ (Section 2.12).
\end{itemize}
\end{warningbox}

\subsection{The ``Clock Reset'' Substitution for Shifted IVPs}

When all of the initial data are stacked at the same point $t = a$ (for example $y(a)$ and $y'(a)$), the cure is to relabel time so that $a$ becomes the new origin. This is a pure change of variable --- it does not alter the differential equation, only the name of the independent variable --- and it places the data exactly where the transform expects them.

\begin{methodbox}[The ``Clock Reset'' Substitution --- Four Steps]
Given an IVP whose conditions are specified at $t = a$:

\smallskip
\textbf{Step 1 --- Reset the clock.} Define the shifted time
\[ \tau = t - a, \qquad\text{equivalently}\qquad t = \tau + a. \]
When $t = a$ the new clock reads $\tau = 0$, so the data move to the origin.

\smallskip
\textbf{Step 2 --- Translate the derivatives (chain rule).} Write $y(t) = y(\tau + a)$ as a function of $\tau$. Since $\dfrac{d\tau}{dt} = \dfrac{d}{dt}(t - a) = 1$, the chain rule gives
\[ \frac{dy}{dt} = \frac{dy}{d\tau}\cdot\frac{d\tau}{dt} = \frac{dy}{d\tau}\cdot 1 = \frac{dy}{d\tau}. \]
Differentiating once more, again with $\frac{d\tau}{dt} = 1$,
\[ \frac{d^2 y}{dt^2} = \frac{d}{dt}\!\left(\frac{dy}{d\tau}\right) = \frac{d}{d\tau}\!\left(\frac{dy}{d\tau}\right)\frac{d\tau}{dt} = \frac{d^2 y}{d\tau^2}. \]
Because the shift factor is exactly $1$, \emph{every derivative is unchanged in form}: $y'(t) = y'(\tau)$ and $y''(t) = y''(\tau)$. The ODE keeps its coefficients verbatim.

\smallskip
\textbf{Step 3 --- Solve entirely in the $\tau$ domain.} The problem now reads with conditions at $\tau = 0$, where $\mathcal{L}\{y'\} = sY - y(0)$ applies cleanly. Transform, solve for $Y(s)$, and invert to get $y(\tau)$.

\smallskip
\textbf{Step 4 --- Slide the clock back.} Undo the substitution with $\tau = t - a$, replacing every $\tau$ in the final answer to recover $y(t)$.
\end{methodbox}

\begin{examplebox}[Worked Example 1: A Shifted Second-Order IVP]
Solve $y'' + y = 0$ with $y\!\left(\tfrac{\pi}{2}\right) = 1$ and $y'\!\left(\tfrac{\pi}{2}\right) = 0$.

\smallskip
\textbf{Step 1 --- Reset the clock.} The data sit at $a = \tfrac{\pi}{2}$, so let
\[ \color{workpurple} \tau = t - \frac{\pi}{2}, \qquad t = \tau + \frac{\pi}{2}. \]
At $t = \tfrac{\pi}{2}$ we have $\tau = 0$, so the conditions become $y(0) = 1$ and $y'(0) = 0$ in the $\tau$ clock.

\smallskip
\textbf{Step 2 --- Translate the ODE.} By the chain rule with $\frac{d\tau}{dt} = 1$, derivatives carry over unchanged, so the equation is identical in form:
\[ \color{workpurple} y''(\tau) + y(\tau) = 0, \qquad y(0) = 1,\quad y'(0) = 0. \]

\textbf{Step 3 --- Transform.} Apply $\mathcal{L}\{y''\} = s^2 Y - s\,y(0) - y'(0)$ with $y(0) = 1$, $y'(0) = 0$:
\[ \color{workpurple} s^2 Y(s) - s(1) - 0 + Y(s) = 0 \quad\Longrightarrow\quad (s^2 + 1)\,Y(s) = s. \]

\textbf{Step 4 --- Solve for $Y(s)$.}
\[ \color{workpurple} Y(s) = \frac{s}{s^2 + 1}. \]

\textbf{Step 5 --- Invert.} Matching $\dfrac{s}{s^2 + 1} \to \cos\tau$:
\[ \color{workpurple} y(\tau) = \cos\tau. \]

\textbf{Step 6 --- Slide the clock back.} Replace $\tau = t - \tfrac{\pi}{2}$:
\[ \color{workpurple} y(t) = \cos\!\left(t - \frac{\pi}{2}\right) = \sin t. \]
The last step uses the identity $\cos\!\left(\theta - \tfrac{\pi}{2}\right) = \sin\theta$. A quick check confirms it: $\sin\!\left(\tfrac{\pi}{2}\right) = 1$ matches $y\!\left(\tfrac{\pi}{2}\right) = 1$, and $y'(t) = \cos t$ gives $\cos\!\left(\tfrac{\pi}{2}\right) = 0$, matching $y'\!\left(\tfrac{\pi}{2}\right) = 0$.
\end{examplebox}

\begin{conceptbox}[Key Takeaway]
For a shifted IVP, the substitution \textcolor{accentorange}{$\tau = t - a$} is a rigid translation of the time axis. Because \textcolor{accentorange}{$\frac{d\tau}{dt} = 1$}, the derivatives --- and therefore the ODE --- are untouched; only the \emph{location} of the data moves, landing at $\tau = 0$ where Laplace can finally use it.
\end{conceptbox}

\subsection{Boundary Value Problems --- Parameterizing the Missing Data}

A boundary value problem (BVP) hands you conditions at \emph{two different} points --- typically one at $t = 0$ and one at a far endpoint $t = b$. The clock-reset trick cannot rescue this case: sliding the axis to put one condition at the origin only pushes the other condition off the origin. We need a different idea.

The transform demands a complete set of data at $t = 0$: for a second-order equation it wants both $y(0)$ \emph{and} $y'(0)$. A BVP usually gives one of them and withholds the other. The resolution is to \textbf{name the missing piece}.

\begin{methodbox}[Parameterizing a Boundary Value Problem]
Given a second-order BVP with one condition at $t = 0$ and one at $t = b$:

\smallskip
\textbf{Step 1 --- Name the unknown zero-data.} Suppose $y(0)$ is given but $y'(0)$ is not. Introduce a symbolic constant
\[ y'(0) = c, \]
treating it as a known value for now. (If instead $y(0)$ is the missing one, name \emph{that} $c$.)

\smallskip
\textbf{Step 2 --- Transform with the symbol carried along.} Apply $\mathcal{L}\{y''\} = s^2 Y - s\,y(0) - y'(0)$, leaving $c$ in place. Solve for $Y(s)$ as an expression that contains $c$.

\smallskip
\textbf{Step 3 --- Invert to a general solution in $c$.} The inverse transform produces $y(t)$ as a family of solutions parameterized by the single unknown $c$.

\smallskip
\textbf{Step 4 --- Apply the far boundary condition.} Substitute the second condition (at $t = b$) into $y(t)$ --- or into $y'(t)$, whichever the condition specifies. This yields one linear equation in the single unknown $c$. Solve it, and back-substitute to get the explicit solution.
\end{methodbox}

\begin{examplebox}[Worked Example 2: A Mixed Boundary Value Problem]
Solve $y'' - y = 0$ with $y(0) = 1$ and $y'(2) = 1$.

\smallskip
\textbf{Step 1 --- Name the missing data.} We know $y(0) = 1$ but not $y'(0)$. Set
\[ \color{workpurple} y'(0) = c. \]

\textbf{Step 2 --- Transform, carrying $c$.} With $\mathcal{L}\{y''\} = s^2 Y - s\,y(0) - y'(0) = s^2 Y - s - c$:
\[ \color{workpurple} s^2 Y(s) - s - c - Y(s) = 0 \quad\Longrightarrow\quad (s^2 - 1)\,Y(s) = s + c. \]

\textbf{Step 3 --- Solve for $Y(s)$ and split.}
\[ \color{workpurple} Y(s) = \frac{s + c}{s^2 - 1} = \frac{s}{s^2 - 1} + \frac{c}{s^2 - 1}. \]

\textbf{Step 4 --- Invert.} Using $\dfrac{s}{s^2 - 1} \to \cosh t$ and $\dfrac{1}{s^2 - 1} \to \sinh t$:
\[ \color{workpurple} y(t) = \cosh t + c\,\sinh t. \]
This is the general solution, still carrying the unknown $c$.

\smallskip
\textbf{Step 5 --- Differentiate for the boundary condition.} The condition is on $y'$, so differentiate, using $\frac{d}{dt}\cosh t = \sinh t$ and $\frac{d}{dt}\sinh t = \cosh t$:
\[ \color{workpurple} y'(t) = \sinh t + c\,\cosh t. \]

\textbf{Step 6 --- Apply $y'(2) = 1$ and solve for $c$.} Substitute $t = 2$:
\[ \color{workpurple} y'(2) = \sinh 2 + c\,\cosh 2 = 1 \quad\Longrightarrow\quad c\,\cosh 2 = 1 - \sinh 2. \]
Isolate $c$ by dividing by $\cosh 2$:
\[ \color{workpurple} c = \frac{1 - \sinh 2}{\cosh 2} \approx \frac{1 - 3.6269}{3.7622} \approx -0.6982. \]

\textbf{Step 7 --- Write the explicit solution.} Back-substitute the value of $c$:
\[ \color{workpurple} y(t) = \cosh t + \frac{1 - \sinh 2}{\cosh 2}\,\sinh t. \]
By construction $y(0) = \cosh 0 = 1$, and $c$ was chosen precisely so that $y'(2) = 1$ --- both boundary conditions are satisfied.
\end{examplebox}

\begin{conceptbox}[Key Takeaway]
A BVP withholds part of the $t = 0$ data the transform requires. \textcolor{accentorange}{Name the missing value $c$}, solve the ODE as a one-parameter family $y(t; c)$, then let the \textcolor{accentorange}{far boundary condition supply the single equation} that fixes $c$. The Laplace machinery runs untouched; the only new work is one linear solve at the end.
\end{conceptbox}

%======================================================================
\clearpage
\phantomsection
\addcontentsline{toc}{section}{Part II --- Practice Set}
\section*{Part II --- Practice Set}

Work each problem before consulting the complete solutions in Part III. Decide first which tool applies: a single shifted data point calls for the clock reset ($\tau = t - a$); split conditions call for parameterizing the missing zero-data with a constant.

\bigskip
\noindent\textbf{Problem 1: First-Order Shifted IVP.}\\[2pt]
Solve
\[ y' + 2y = 0, \qquad y(3) = 4. \]
\textit{Hint: the data sit at $t = 3$. Reset the clock with $\tau = t - 3$.}

\bigskip
\noindent\textbf{Problem 2: Second-Order Shifted IVP.}\\[2pt]
Solve
\[ y'' + 4y = 0, \qquad y(\pi) = 2,\quad y'(\pi) = -4. \]
\textit{Hint: both conditions sit at $t = \pi$. Reset with $\tau = t - \pi$, then simplify using periodicity at the end.}

\bigskip
\noindent\textbf{Problem 3: Mixed Boundary Value Problem.}\\[2pt]
Solve
\[ y'' + 9y = 0, \qquad y(0) = 2,\quad y\!\left(\tfrac{\pi}{6}\right) = 1. \]
\textit{Hint: $y(0)$ is given but $y'(0)$ is not. Let $y'(0) = c$ and use the far condition to solve for it.}

%======================================================================
\clearpage
\phantomsection
\addcontentsline{toc}{section}{Part III --- Complete Solutions}
\section*{Part III --- Complete Solutions}

\noindent\textbf{Solution 1: First-Order Shifted IVP.}\\[2pt]
The single data point lives at $t = 3$, so we reset the clock.

\smallskip
\textbf{Step 1 --- Reset the clock.} Let
\[ \color{workpurple} \tau = t - 3, \qquad t = \tau + 3. \]
At $t = 3$ the new clock reads $\tau = 0$, so the condition becomes $y(0) = 4$.

\smallskip
\textbf{Step 2 --- Translate the ODE.} Since $\frac{d\tau}{dt} = 1$, the chain rule gives $y'(t) = y'(\tau)$, so the equation is unchanged in form:
\[ \color{workpurple} y'(\tau) + 2y(\tau) = 0, \qquad y(0) = 4. \]

\textbf{Step 3 --- Transform.} Apply $\mathcal{L}\{y'\} = sY - y(0)$ with $y(0) = 4$:
\[ \color{workpurple} sY(s) - 4 + 2Y(s) = 0 \quad\Longrightarrow\quad (s + 2)\,Y(s) = 4. \]

\textbf{Step 4 --- Solve for $Y(s)$.}
\[ \color{workpurple} Y(s) = \frac{4}{s + 2}. \]

\textbf{Step 5 --- Invert.} Using $\dfrac{1}{s + 2} \to e^{-2\tau}$:
\[ \color{workpurple} y(\tau) = 4e^{-2\tau}. \]

\textbf{Step 6 --- Slide the clock back.} Replace $\tau = t - 3$:
\[ \color{workpurple} y(t) = 4e^{-2(t - 3)}. \]
Check: $y(3) = 4e^{0} = 4$, as required.

\bigskip
\noindent\textbf{Solution 2: Second-Order Shifted IVP.}\\[2pt]
Both conditions sit at $t = \pi$, so a single clock reset places them at the origin.

\smallskip
\textbf{Step 1 --- Reset the clock.} Let
\[ \color{workpurple} \tau = t - \pi, \qquad t = \tau + \pi. \]
The conditions become $y(0) = 2$ and $y'(0) = -4$ in the $\tau$ clock.

\smallskip
\textbf{Step 2 --- Translate the ODE.} With $\frac{d\tau}{dt} = 1$, both derivatives carry over unchanged ($y'(t) = y'(\tau)$, $y''(t) = y''(\tau)$), so
\[ \color{workpurple} y''(\tau) + 4y(\tau) = 0, \qquad y(0) = 2,\quad y'(0) = -4. \]

\textbf{Step 3 --- Transform.} Apply $\mathcal{L}\{y''\} = s^2 Y - s\,y(0) - y'(0)$ with $y(0) = 2$, $y'(0) = -4$:
\[ \color{workpurple} s^2 Y(s) - 2s - (-4) + 4Y(s) = 0 \quad\Longrightarrow\quad (s^2 + 4)\,Y(s) = 2s - 4. \]

\textbf{Step 4 --- Solve for $Y(s)$ and split.}
\[ \color{workpurple} Y(s) = \frac{2s - 4}{s^2 + 4} = 2\cdot\frac{s}{s^2 + 4} - 2\cdot\frac{2}{s^2 + 4}. \]
The second term is written with a $2$ in the numerator to match the standard form $\dfrac{2}{s^2 + 2^2} \to \sin(2\tau)$.

\smallskip
\textbf{Step 5 --- Invert.} Using $\dfrac{s}{s^2 + 4} \to \cos(2\tau)$ and $\dfrac{2}{s^2 + 4} \to \sin(2\tau)$:
\[ \color{workpurple} y(\tau) = 2\cos(2\tau) - 2\sin(2\tau). \]

\textbf{Step 6 --- Slide the clock back.} Replace $\tau = t - \pi$:
\[ \color{workpurple} y(t) = 2\cos\!\big(2(t - \pi)\big) - 2\sin\!\big(2(t - \pi)\big). \]
Now simplify using $2(t - \pi) = 2t - 2\pi$ and the $2\pi$-periodicity of sine and cosine, so $\cos(2t - 2\pi) = \cos 2t$ and $\sin(2t - 2\pi) = \sin 2t$:
\[ \color{workpurple} y(t) = 2\cos(2t) - 2\sin(2t). \]
Check: $y(\pi) = 2\cos(2\pi) - 2\sin(2\pi) = 2$, and $y'(t) = -4\sin(2t) - 4\cos(2t)$ gives $y'(\pi) = -4\cos(2\pi) = -4$. Both conditions hold.

\bigskip
\noindent\textbf{Solution 3: Mixed Boundary Value Problem.}\\[2pt]
The conditions are split between $t = 0$ and $t = \tfrac{\pi}{6}$, so we parameterize the missing initial slope.

\smallskip
\textbf{Step 1 --- Name the missing data.} We have $y(0) = 2$ but not $y'(0)$. Set
\[ \color{workpurple} y'(0) = c. \]

\textbf{Step 2 --- Transform, carrying $c$.} With $\mathcal{L}\{y''\} = s^2 Y - s\,y(0) - y'(0) = s^2 Y - 2s - c$:
\[ \color{workpurple} s^2 Y(s) - 2s - c + 9Y(s) = 0 \quad\Longrightarrow\quad (s^2 + 9)\,Y(s) = 2s + c. \]

\textbf{Step 3 --- Solve for $Y(s)$ and split.}
\[ \color{workpurple} Y(s) = \frac{2s + c}{s^2 + 9} = 2\cdot\frac{s}{s^2 + 9} + \frac{c}{3}\cdot\frac{3}{s^2 + 9}. \]
The factor of $\tfrac{1}{3}$ appears so the last term matches $\dfrac{3}{s^2 + 3^2} \to \sin(3t)$.

\smallskip
\textbf{Step 4 --- Invert.} Using $\dfrac{s}{s^2 + 9} \to \cos(3t)$ and $\dfrac{3}{s^2 + 9} \to \sin(3t)$:
\[ \color{workpurple} y(t) = 2\cos(3t) + \frac{c}{3}\,\sin(3t). \]

\textbf{Step 5 --- Apply the far boundary condition.} Substitute $t = \tfrac{\pi}{6}$, noting $3\cdot\tfrac{\pi}{6} = \tfrac{\pi}{2}$, so $\cos\!\left(\tfrac{\pi}{2}\right) = 0$ and $\sin\!\left(\tfrac{\pi}{2}\right) = 1$:
\[ \color{workpurple} y\!\left(\tfrac{\pi}{6}\right) = 2\cos\!\left(\tfrac{\pi}{2}\right) + \frac{c}{3}\sin\!\left(\tfrac{\pi}{2}\right) = 2(0) + \frac{c}{3}(1) = \frac{c}{3}. \]
Set this equal to the given value $1$:
\[ \color{workpurple} \frac{c}{3} = 1 \quad\Longrightarrow\quad c = 3. \]

\textbf{Step 6 --- Write the explicit solution.} Back-substitute $c = 3$, so $\tfrac{c}{3} = 1$:
\[ \color{workpurple} y(t) = 2\cos(3t) + \sin(3t). \]
Check: $y(0) = 2\cos 0 + \sin 0 = 2$, and $y\!\left(\tfrac{\pi}{6}\right) = 2(0) + 1 = 1$. Both boundary conditions are satisfied.

\begin{conceptbox}[Closing Takeaway]
Two questions decide the route every time. \emph{Where do the conditions live?} If they are stacked at one point $t = a$, \textcolor{accentorange}{reset the clock with $\tau = t - a$} so they land at the origin. If they are \emph{split} between $t = 0$ and a far point, \textcolor{accentorange}{name the missing zero-data $c$} and let the far condition solve for it. In both cases the Laplace transform itself never changes --- only the bookkeeping that delivers usable data to $t = 0$.
\end{conceptbox}

\end{document}
